\chapter{Diseño de la solución}\label{cap:diseno}

\section{Introducción}
El diseño persigue un equilibrio explícito entre tres objetivos: capacidad analítica,
reproducibilidad local y claridad para defensa académica. En lugar de una arquitectura
sobredimensionada, se adopta un diseño modular con separación entre datos, notebooks y
documentación, lo que permite ampliar componentes sin rehacer el núcleo del flujo
analítico.

\section{Arquitectura de alto nivel}
La solución se organiza en cuatro capas:

\begin{itemize}
	\item \textbf{Capa de ejecución}: contenedores \texttt{zeppelin},
	\texttt{spark-master} y \texttt{spark-worker-1} en Docker Compose.
	\item \textbf{Capa de datos}: directorio \texttt{data/} con el Parquet de
	entrada, CSVs de indicadores exportados, shapefile y GeoJSON de zonas.
	\item \textbf{Capa de procesamiento}: notebooks de Zeppelin con PySpark y
	Spark SQL para carga, limpieza, análisis y exportación.
	\item \textbf{Capa de visualización}: gráficas nativas de Zeppelin para
	análisis tabular y mapas coropléticos con Folium para análisis geográfico.
\end{itemize}

El uso de Spark en modo distribuido local permite conservar la semántica de un
procesamiento escalable, aun ejecutando en infraestructura docente.

\section{Diseño de almacenamiento y rutas}
El diseño físico de datos responde a un flujo directo:

\begin{itemize}
	\item \texttt{data/yellow\_tripdata\_2025-01.parquet}: fichero de entrada
	con los registros Yellow Taxi de enero de 2025.
	\item \texttt{data/taxi\_zones.*}: shapefile oficial de zonas TLC (5 ficheros).
	\item \texttt{data/nyc\_taxi\_zones.geojson}: GeoJSON generado a partir del
	shapefile mediante el script auxiliar \texttt{convert.py}.
	\item \texttt{data/*.csv}: indicadores por zona exportados desde Spark
	(14 ficheros: viajes, tarifa, propina, distancia, duración, ingresos y horas
	pico/valle, cada uno por origen y por destino).
\end{itemize}

Esta estructura plana es intencionada: todos los datos residen en un único directorio
montado como volumen Docker (\texttt{/data}), accesible tanto desde Spark como desde
Python.

\section{Diseño del flujo analítico}
El flujo se implementa íntegramente en dos notebooks de Zeppelin, sin scripts
externos de ETL:

\begin{enumerate}
	\item \textbf{Notebook principal}
	(\texttt{Análisis Taxis Nueva York.zpln}): concentra la carga del Parquet,
	la limpieza, el enriquecimiento y todas las consultas analíticas. Está
	organizado en bloques secuenciales que procesan los datos paso a paso.
	\item \textbf{Notebook de mapas}
	(\texttt{Análisis Taxis Nueva York Mapas.zpln}): lee los CSVs de indicadores
	exportados por el notebook principal y genera mapas coropléticos interactivos
	con Folium.
\end{enumerate}

Se incluye además una versión del notebook de mapas ya ejecutada
(\texttt{Análisis Taxis Nueva York Mapas (Ejecutados).zpln}), que permite revisar
las visualizaciones sin necesidad de volver a ejecutar el código.

En la fase de procesamiento se definen métricas derivadas clave:

\[
\text{trip\_duration\_minutes} =
\frac{\text{unix}(\texttt{dropoff}) - \text{unix}(\texttt{pickup})}{60}
\]

\[
\text{average\_speed\_mph} =
\frac{\texttt{trip\_distance}}{\text{trip\_duration\_minutes} / 60}
\]

Y se aplican restricciones de limpieza para eliminar registros no plausibles
(duraciones no positivas, velocidades extremas, importes negativos y nulos críticos).

\section{Diseño del notebook principal}
El notebook principal se diseña como un flujo integrado con párrafos Spark y
Spark SQL:

\begin{enumerate}
	\item Carga del Parquet de referencia y validación básica de esquema.
	\item Limpieza avanzada con filtros de dominio (rango de
	\texttt{LocationID}, pasajeros mínimos y coherencia temporal).
	\item Enriquecimiento geográfico mediante GeoJSON de zonas TLC.
	\item Traducción de códigos a etiquetas legibles (vendor, pago, tarifa, flag).
	\item Ejecución de consultas analíticas por hora, franja, tipo de viaje,
	pasajeros y balance entre recogidas y destinos.
	\item Exportación de indicadores agregados por zona a CSV.
	\item Generación de visualizaciones comparativas directamente en Zeppelin.
\end{enumerate}

Este diseño permite que el notebook sea simultáneamente entorno de análisis,
artefacto de demostración y evidencia de ejecución del proyecto.

\section{Diseño de consultas e indicadores}
Los indicadores extraídos se agrupan en cuatro vistas principales:

\begin{itemize}
\item \textbf{Comportamiento temporal}: viajes por hora, tarifa media y duración
media, además de comparativa de 1 de enero frente al resto del mes.
\item \textbf{Comparativas operativas}: diferencia entre viajes estándar y viajes
de aeropuerto en volumen, distancia, duración, tarifa y propina.
\item \textbf{Patrones geográficos}: métricas agregadas por distrito de recogida
(distancia, velocidad y coste por milla).
\item \textbf{Balance territorial}: contraste recogidas vs destinos por distrito
y su diferencia neta.
\end{itemize}

Los indicadores por zona se exportan a CSV con la siguiente estructura:
\begin{itemize}
\item \textbf{Por zona de origen (pickup)}: viajes, tarifa media, propina media,
distancia media, duración media, ingresos totales y hora pico/valle.
\item \textbf{Por zona de destino (dropoff)}: los mismos indicadores agregados por
\texttt{DOLocationID}.
\end{itemize}

Estos CSVs alimentan directamente el notebook de mapas, cerrando el ciclo
análisis-visualización.

\section{Decisiones de diseño y trade-offs}
Las decisiones más relevantes fueron:

\begin{itemize}
	\item \textbf{Parquet como fuente principal}: mejora de I/O y lectura selectiva
	por columnas, adecuado para cargas analíticas con Spark.
	\item \textbf{CSV como formato de exportación}: los indicadores por zona se
	exportan a CSV para su consumo por Folium/Pandas, que requiere formatos
	tabulares simples. No se usa CSV como formato de almacenamiento analítico.
	\item \textbf{Zeppelin como interfaz única}: facilita defensa al unir código,
	salida, gráficas y explicación en el mismo documento.
	\item \textbf{Dos notebooks separados}: el análisis principal se mantiene
	independiente de la visualización geográfica, permitiendo ejecutar uno sin
	el otro y reduciendo el acoplamiento.
	\item \textbf{Consultas SQL centradas en interpretación}: se priorizan métricas
	que aportan lectura de movilidad real frente a tablas excesivamente técnicas.
\end{itemize}

Como trade-off, se renuncia a complejidad de producción (autoscaling,
alta disponibilidad y orquestación externa) para mantener foco académico.

\section{Conclusiones}
El diseño implementado es coherente con el alcance del proyecto de asignatura y con
la naturaleza del dataset. La separación en dos notebooks, junto con la exportación
de indicadores por zona, proporciona una base sólida para análisis reproducible hoy y
para evolución a despliegues más escalables en el futuro.